\documentclass[troisieme]{fiche}
\metadonnees{19}{Ce que voient les étoiles}{Bloc D — Rapporter et supposer}

\begin{document}
\entete

\objectifs{%
\textbf{Langue} : l'hypothèse irréelle, \emph{if} + prétérit ; \emph{I wish}.\par
\textbf{Texte} : J.\,M. Barrie, \emph{Peter Pan} (1911), ch. \textsc{ii}.\par
\textbf{Durée} : 1 h — exercices 20 min, texte et questions 25 min, version 15 min.}

%% ===================================================================
\exercice[12 min]{L'irréel du présent}

\rappel{\textbf{Rappel.} \emph{If} + prétérit dans la condition, \emph{would} dans la
conséquence. Avec \emph{be}, on emploie \emph{were} à toutes les personnes.}

\consigne{Mettez le verbe entre parenthèses à la forme qui convient.}

\begin{anglais}
\begin{enumerate}[label=\arabic*.,itemsep=2.5mm,leftmargin=*]
  \item If I \emph{(be)} \trou{were} a bird, I \emph{(fly)} \trou{would fly} to Egypt.
  \item I wish I \emph{(have)} \trou{had} a dog like Nana.
  \item If the window \emph{(be)} \trou{were} open, the boy \emph{(can)} \trou{could} come in.
  \item She wishes she \emph{(not go)} \trou{were not going} to the party tonight.
  \item If Mr Darling \emph{(listen)} \trou{listened} to his wife, nothing \emph{(happen)}
    \trou{would happen}.
  \item I wish you \emph{(can)} \trou{could} hear the fairy language.
  \item If you \emph{(hear)} \trou{heard} it, you \emph{(know)} \trou{would know} that you had
    heard it before.
  \item What \emph{(you do)} \trou{would you do} if a boy \emph{(come)} \trou{came} through
    your window?
  \item I wish the stars \emph{(speak)} \trou{spoke} instead of winking.
  \item If Nana \emph{(not be)} \trou{were not} tied up, she \emph{(bark)} \trou{would bark}.
\end{enumerate}
\end{anglais}

\note[Un prétérit qui n'est pas un passé]{\emph{If I were a bird} ne parle pas d'hier mais
d'aujourd'hui : le prétérit y marque l'écart avec le réel, non le temps. D'où \emph{were} à
toutes les personnes — \emph{if I were}, \emph{if he were} — forme qu'on ne rencontre que
là.}

\note[\emph{I wish}]{Suivi du prétérit, il exprime un regret sur le présent : \emph{I wish I
had a dog} = je n'en ai pas (2, 9). Avec un verbe de capacité, \emph{could} (6). Et jamais de
\emph{would} après \emph{I wish} pour soi-même : \emph{I wish I were}, non \emph{I wish I
would be}.}

\note[Jamais \emph{would} après \emph{if}]{Comme après \emph{if} on ne met pas \emph{will},
on n'y met pas \emph{would} : \emph{if he listened}, et non \emph{if he would listen}.}

%% ===================================================================
\exercice[8 min]{Souhaiter, supposer}

\rappel{\textbf{Rappel.} \og Si seulement \fg{} : \emph{if only} + prétérit. \og J'aimerais
que tu\ldots \fg{} : \emph{I wish you would\ldots}}

\consigne{Traduisez en anglais.}

\begin{anglais}
\begin{enumerate}[label=\arabic*.,itemsep=3mm,leftmargin=*]
  \item Si j'étais toi, je ne sortirais pas.
    \lignes{1}
    \corr{If I were you, I wouldn't go out.}
  \item J'aimerais savoir voler.
    \lignes{1}
    \corr{I wish I could fly.}
  \item Si elle savait la vérité, elle pleurerait.
    \lignes{1}
    \corr{If she knew the truth, she would cry.}
  \item Si seulement il était là !
    \lignes{1}
    \corr{If only he were here!}
  \item J'aimerais que tu restes.
    \lignes{1}
    \corr{I wish you would stay.}
  \item Que ferais-tu si tu voyais une fée ?
    \lignes{1}
    \corr{What would you do if you saw a fairy?}
\end{enumerate}
\end{anglais}

\note[Deux \emph{wish}]{Pour soi-même, \emph{I wish} + prétérit (2). Pour ce que fait un
autre, et qu'on voudrait voir changer, \emph{I wish} + \emph{would} (5) : \emph{I wish you
would stay} sous-entend que l'autre s'apprête à partir.}

%% ===================================================================
\newpage
\section{Texte}

\noindent\small\textit{Chez les Darling, à Londres. La chienne Nana sert de nourrice aux trois
enfants. Ce soir-là, vexé, Mr Darling vient de l'attacher dans la cour avant de partir dîner
chez des voisins.}\normalsize

\begin{texte}
In the meantime Mrs. Darling had put the children to bed in \gl[unwonted]{unwonted}{inhabituel}
silence and lit their night-lights. They could hear Nana barking, and John
\gl[to whimper]{whimpered}{gémir}, ``It is because he is chaining her up in the yard,'' but
Wendy was wiser.

``That is not Nana's unhappy bark,'' she said, little \gl[to guess]{guessing}{deviner,
soupçonner} what was about to happen; ``that is her bark when she \gl[to smell]{smells}{sentir,
flairer} danger.''

Danger!

``Are you sure, Wendy?''

``Oh, yes.''

Mrs. Darling \gl[to quiver]{quivered}{frémir} and went to the window. It was
\gl[securely]{securely}{solidement} \gl[to fasten]{fastened}{fermer, attacher}. She looked
out, and the night was \gl[to pepper]{peppered}{parsemer} with stars. They were crowding round
the house, as if curious to see what was to take place there, but she did not notice this, nor
that one or two of the smaller ones \gl[to wink]{winked}{faire un clin d'œil} at her. Yet a
\gl[nameless]{nameless}{sans nom} fear \gl[to clutch]{clutched}{agripper} at her heart and made
her cry, ``Oh, how I wish that I wasn't going to a party to-night!''

Even Michael, already half asleep, knew that she was
\gl[to perturb]{perturbed}{troubler}, and he asked, ``Can anything harm us, mother, after the
night-lights are lit?''

``Nothing, precious,'' she said; ``they are the eyes a mother leaves behind her to guard her
children.''

She went from bed to bed singing \gl[an enchantment]{enchantments}{un charme, une incantation}
over them, and little Michael \gl[to fling, flung]{flung}{jeter} his arms round her.
``Mother,'' he cried, ``I'm glad of you.'' They were the last words she was to hear from him
for a long time.

No. 27 was only a few yards distant, but there had been a slight fall of snow, and Father and
Mother Darling picked their way over it \gl[deftly]{deftly}{adroitement} not to
\gl[to soil]{soil}{salir} their shoes. They were already the only persons in the street, and
all the stars were watching them.
\end{texte}
\source{J.\,M. Barrie, \emph{Peter Pan}, Londres, Hodder \& Stoughton, 1911,
ch.~\textsc{ii}.}

\partiel[15 min]{Questions}
\consigne{Répondez aux deux premières questions en français, aux deux dernières en anglais.}

\begin{enumerate}[label=\arabic*.,itemsep=2.5mm,leftmargin=*]
  \item Pourquoi Nana aboie-t-elle, selon John, et selon Wendy ?
    \lignes{2}
    \corr{Pour John, parce qu'on est en train de l'attacher dans la cour. Wendy, elle,
    reconnaît un autre aboiement : celui que Nana pousse quand elle flaire un danger.}
  \item Que répond Mrs Darling à Michael au sujet des veilleuses ?
    \lignes{2}
    \corr{Que rien ne peut leur arriver : les veilleuses sont les yeux qu'une mère laisse
    derrière elle pour garder ses enfants.}
  \item \en{``They were the last words she was to hear from him for a long time.'' What does
    the narrator do here that none of the characters can?}
    \lignes{3}
    \corr{He steps out of the evening and tells us what is coming, before the mother knows it.
    The reader is made to see the goodbye that Michael and his mother do not see. Everything in
    the scene is already a warning — the barking, the stars, the fastened window — and only
    the narrator can say so.}
  \item \en{Find the sentence that begins \emph{Oh, how I wish\ldots} What tense follows
    \emph{wish}, and what does the sentence tell us about the party?}
    \lignes{3}
    \corr{\emph{Oh, how I wish that I wasn't going to a party to-night!} Un prétérit
    (\emph{wasn't going}) après \emph{wish} : le souhait porte sur maintenant, et il est déjà
    perdu — elle y va. La forme dit à elle seule qu'il est trop tard.}
\end{enumerate}

\note[Ce que la langue laisse passer]{\emph{I wish that I wasn't going} : à l'écrit soigné,
\emph{were} conviendrait mieux que \emph{wasn't}. Barrie fait parler une mère inquiète, pas un
manuel : l'anglais parlé emploie couramment \emph{was} dans ce tour.}

%% ===================================================================
\section{Version}
\consigne{Traduisez le passage en français. Il suit immédiatement le texte.}

\begin{version}
Stars are beautiful, but they may not take an active part in anything, they must just look on
for ever. It is a punishment put on them for something they did so long ago that no star now
knows what it was. So the older ones have become \gl[glassy-eyed]{glassy-eyed}{au regard
vitreux} and \gl[seldom]{seldom}{rarement} speak (winking is the star language), but the
little ones still wonder. They are not really friendly to Peter, who had a
\gl[mischievous]{mischievous}{espiègle, malicieux} way of stealing up behind them and trying to
blow them out; but they are so fond of fun that they were on his side to-night, and anxious to
get the grown-ups out of the way. So as soon as the door of 27 closed on Mr. and Mrs. Darling
there was a \gl[a commotion]{commotion}{une agitation} in the
\gl[the firmament]{firmament}{le firmament}, and the smallest of all the stars in the Milky
Way screamed out: ``Now, Peter!''
\end{version}
\source{\emph{Ibid.}}

\lignes{16}

\corr{\textbf{Proposition de traduction.} Les étoiles sont belles, mais il ne leur est pas
permis de prendre part à quoi que ce soit ; elles doivent seulement regarder, à jamais. C'est
une punition qu'on leur a infligée pour quelque chose qu'elles ont fait il y a si longtemps
qu'aucune étoile ne sait plus ce que c'était. Aussi les plus vieilles sont-elles devenues
vitreuses et parlent-elles rarement (le clin d'œil est la langue des étoiles), mais les
petites s'étonnent encore. Elles n'aiment pas beaucoup Peter, qui avait la fâcheuse habitude
de se glisser derrière elles pour les souffler ; mais elles aiment tant s'amuser qu'elles
étaient de son côté ce soir-là, et pressées d'écarter les grandes personnes. Aussi, dès que la
porte du 27 se fut refermée sur M. et Mme Darling, il y eut une agitation dans le firmament, et
la plus petite de toutes les étoiles de la Voie lactée s'écria : \og Maintenant, Peter ! \fg}

\note[Choix de traduction 1]{\emph{they may not take an active part} : \emph{may not} est ici
l'interdiction, non l'éventualité. \og Il ne leur est pas permis \fg{} ; \og elles ne peuvent
peut-être pas \fg{} serait un contresens.}

\note[Choix de traduction 2]{\emph{a punishment put on them} : participe passé seul, sans
relatif — \emph{which was put on them}. Le français rétablit \og qu'on leur a infligée \fg.}

\note[Choix de traduction 3]{\emph{so\ldots that} et \emph{so fond of fun that} : la
conséquence. Le français dit \og si\ldots que \fg{} ou \og tant\ldots que \fg{} selon qu'il
s'agit d'un adjectif ou d'un verbe.}

%% ===================================================================
\partiel{Lexique à mémoriser}
\begin{multicols}{2}\small
\begin{itemize}[label=--,itemsep=0.8mm,leftmargin=*]
  \item \textbf{to bark} — aboyer
  \item \textbf{a yard} — une cour \textit{(fiche 14)}
  \item \textbf{to smell} — sentir ; smelt, smelt
  \item \textbf{danger} — le danger ; \textit{dangerous} dangereux
  \item \textbf{to fasten} — fermer, attacher
  \item \textbf{a star} — une étoile
  \item \textbf{to wink} — faire un clin d'œil
  \item \textbf{fear} — la peur ; \textit{to be afraid} avoir peur
  \item \textbf{to harm} — faire du mal à
  \item \textbf{to guard} — garder, protéger
  \item \textbf{a punishment} — une punition
  \item \textbf{seldom} — rarement
\end{itemize}
\end{multicols}

\end{document}
